\section{Evaluations}


This section begins with the post-training evaluation of Kimi-K2-Instruct, followed by a brief overview of the capabilities of Kimi-K2-Base. We conclude with a comprehensive safety evaluation.


\subsection{Post-training Evaluations}



\subsubsection{Evaluation Settings}


\paragraph{Benchmarks}
We assess Kimi-K2-Instruct across different areas.
For coding, we adopt LiveCodeBench v6~\citep{jain2024livecodebench}(questions from August 2024 to May 2025), OJBench~\citep{wang2025ojbenchcompetitionlevelcode}, MultiPL-E~\citep{10103177}, SWE-bench Verified~\citep{jimenez2024swebench,yang2025swesmith}, TerminalBench~\citep{tbench_2025}, Multi-SWE-bench~\citep{zan2025multi}, SWE-Lancer~\citep{miserendino2025swelancer}, PaperBench~\citep{starace2025paperbench}, and Aider-Polyglot~\citep{aider}. 
For tool use tasks, we evaluate performance on $\tau^2$-Bench~\citep{barres2025tau2} and AceBench~\citep{chen2025acebench}, which emphasize multi-turn tool-calling capabilities.
In reasoning, we include a wide range of mathematical, science and logical tasks: AIME 2024/2025, MATH-500, HMMT 2025, CNMO 2024, PolyMath-en, ZebraLogic~\citep{lin2025zebralogicscalinglimitsllms}, AutoLogi~\citep{zhu2025autologiautomatedgenerationlogic}, GPQA-Diamond~\citep{rein2024gpqa}, SuperGPQA~\citep{du2025supergpqa}, and Humanity's Last Exam (Text-Only)~\citep{phan2025humanitysexam}.
We benchmark the long-context capabilities on:  MRCR\footnote{\url{https://huggingface.co/datasets/openai/mrcr}} for long-context retrieval, and  DROP~\citep{DBLP:journals/corr/abs-1903-00161}, FRAMES~\citep{krishna2025factfetchreasonunified}  and  LongBench v2~\citep{bai2025longbenchv2deeperunderstanding} for  long-context reasoning.
For factuality, we evaluate FACTS Grounding~\citep{jacovi2025factsgroundingleaderboardbenchmarking}, the Vectara Hallucination Leaderboard~\citep{hhem-2.1-open}, and FaithJudge~\citep{tamber2025benchmarking}.
Finally, general capabilities are assessed using MMLU~\citep{hendrycks2021measuringmassivemultitasklanguage}, MMLU-Redux~\citep{gema2024we}, MMLU-Pro~\citep{wang2024mmluprorobustchallengingmultitask}, IFEval~\citep{Zhou2023InstructionFollowingEF}, Multi-Challenge~\citep{sirdeshmukh2025multichallengerealisticmultiturnconversation}, SimpleQA~\citep{wei2024measuring}, and LiveBench~\citep{livebench} (as of 2024-11-25).


\paragraph{Baselines}{
We benchmark against both open-source and proprietary frontier models, ensuring every candidate is evaluated under its non-thinking configuration to eliminate additional gains from test-time compute.
Open-source baselines: DeepSeek-V3-0324 and Qwen3-235B-A22B, with the latter run in the vendor-recommended no-thinking regime.
Proprietary baselines: Claude Sonnet 4, Claude Opus 4, GPT-4.1, and Gemini 2.5 Flash Preview (2025-05-20). Each invoked in its respective non-thinking mode via official APIs under unified temperature and top-p settings.

\textbf{Evaluation Configurations} All runs query models in their non-thinking mode. Output token length is capped at 8192 tokens everywhere except SWE-bench Verified (Agentless), which is raised to 16384. 
For benchmarks with high per-question variance, we adopt repeated sampling $k$ times and average the results to obtain stable scores, denoted as Avg@k.
For long-context tasks, we set the context window size to 128K tokens during evaluation, truncating any input that exceeds this limit to fit within the window.
SWE-bench Verified is evaluated in two modes: Agentless Coding via Single Patch without Test (Acc) and Agentic Coding via bash/editor tools under both Single Attempt (Acc) and Multiple Attempts (Acc) using best-of-N selection with an internal verifier; SWE-bench Multilingual is tested only in the single-attempt agentic setting. Some data points have been omitted due to prohibitively expensive evaluation costs.
}

\newcolumntype{P}[1]{>{\centering\arraybackslash}p{#1}}

\begin{table}[htbp] % 使用 [htbp] 提供更灵活的浮动位置
\centering
\footnotesize % 使用 footnotesize，比 scriptsize 稍大，更易读
\setlength{\tabcolsep}{4pt} % 减少列间距以节省空间

\caption{Performance comparison of Kimi-K2-Instruct against leading open-source and proprietary models across diverse tasks. \textbf{Bold} denotes the global SOTA; \underline{\textbf{underlined bold}} indicates the best open-source result. Data points marked with * are taken directly from the model's technical report or blog.}
\label{tab:instruct_eval} % 添加一个标签，方便在文中引用
\begin{tabular}{@{}l *{7}{P{1.4cm}}@{}} % *{7}{P{1.6cm}} 表示重复 P{1.6cm} 7次
\toprule
& \multicolumn{3}{c}{\textbf{Open Source}} & \multicolumn{4}{c}{\textbf{Proprietary}} \\ % 调整了 & 数量
\cmidrule(lr){2-4} \cmidrule(lr){5-8} % 调整了 cmidrule 的列范围
\textbf{Benchmark} & \textbf{Kimi-K2-Instruct} & \textbf{DeepSeek-V3-0324} & \textbf{Qwen3-235B-A22B} & \textbf{Claude Sonnet 4} & \textbf{Claude Opus 4} & \textbf{GPT-4.1} & \textbf{Gemini 2.5 Flash} \\
\midrule
\multicolumn{8}{@{}l}{\textbf{Coding Tasks}} \\ % 调整为 8 列
\midrule
LiveCodeBench~v6 (Pass@1) & \textbf{53.7} & 46.9 & 37.0 & 48.5 & 47.4 & 44.7 & 44.7 \\
OJBench (Pass@1) & \textbf{27.1} & 24.0 & 11.3 & 15.3 & 19.6 & 19.5 & 19.5 \\
MultiPL-E (Pass@1) & \textbf{\underline{85.7}} & 83.1 & 78.2 & 88.6 & \textbf{89.6} & 86.7 & 85.6 \\
\makecell[l]{SWE-bench~Verified\\\textit{Agentless-Single-Patch} (Pass@1)} & \textbf{\underline{51.8}} & 36.6 & 39.4 & 50.2 & \textbf{53.0} & 40.8 & 32.6 \\
\makecell[l]{SWE-bench~Verified\\\textit{Agentic-Single-Attempt} (Pass@1)} & \textbf{\underline{65.8}} & 38.8 & 34.4 & \textbf{72.7*} & 72.5* & 54.6 & --- \\
\makecell[l]{SWE-bench~Verified\\\textit{Agentic-Multi-Attempt} (Pass@1)} & \textbf{\underline{71.6}} & --- & --- & \textbf{80.2*} & 79.4* & --- & --- \\
\makecell[l]{SWE-bench Multilingual (Pass@1)} & \textbf{\underline{47.3}} & 25.8 & 20.9 & \textbf{51.0} & --- & 31.5 & --- \\
\makecell[l]{Multi-SWE-bench (Pass@1)} & \textbf{\underline{18.3}} & 8.0 & 9.0 & \textbf{29.2}  & --- & 11.7 & 14.0 \\
\makecell[l]{SWE-Lancer (Pass@1)} & \textbf{\underline{39.1}} & 30.5 & 24.1 & \textbf{40.8}  & --- & 23.0 & 38.5 \\
\makecell[l]{Paper Bench~\textit{Code-Dev} (Acc.)} & \textbf{\underline{27.8}} & 12.2 & 13.2 & \textbf{43.3}  & --- & 29.9 & 5.7 \\

Terminal Bench \textit{In-House} (Acc.) & \textbf{\underline{30.0}} & --- & --- & 35.5 & \textbf{43.2} & 8.3 & --- \\
Terminal Bench \textit{Terminus} (Acc.) & \textbf{\underline{25.0}} & 16.3 & 6.6 & --- & --- & \textbf{30.3} & 16.8 \\
Aider-Polyglot (Acc.) & 60.0 & 55.1 & \textbf{\underline{61.8}} & 56.4 & \textbf{70.7} & 52.4 & 44.0 \\
\midrule
\multicolumn{8}{@{}l}{\textbf{Tool Use Tasks}} \\ % 调整为 8 列
\midrule
Tau2 retail (Avg@4) & \textbf{\underline{70.6}} & 69.1 & 57.0 & 75.0 & \textbf{81.8} & 74.8 & 64.3 \\
Tau2 airline (Avg@4) & \textbf{\underline{56.5}} & 39.0 & 26.5 & 55.5 & \textbf{60.0} & 54.5 & 42.5 \\
Tau2 telecom (Avg@4) & \textbf{65.8} & 32.5 & 22.1 & 45.2 & 57.0 & 38.6 & 16.9 \\
AceBench (Acc.) & \textbf{\underline{76.5}} & 72.7 & 70.5 & 76.2 & 75.6 & \textbf{80.1} & 74.5 \\
\midrule
\multicolumn{8}{@{}l}{\textbf{Math \& STEM Tasks}} \\ % 调整为 8 列
\midrule
AIME 2024 (Avg@64) & \textbf{69.6} & 59.4* & 40.1* & 43.4 & 48.2 & 46.5 & 61.3 \\
AIME 2025 (Avg@64) & \textbf{49.5} & 46.7 & 24.7* & 33.1* & 33.9* & 37.0 & 46.6 \\
MATH-500 (Acc.) & \textbf{97.4} & 94.0* & 91.2* & 94.0 & 94.4 & 92.4 & 95.4 \\
HMMT 2025 (Avg@32) & \textbf{38.8} & 27.5 & 11.9 & 15.9 & 15.9 & 19.4 & 34.7 \\
CNMO 2024 (Avg@16) & 74.3 & \textbf{\underline{74.7}} & 48.6 & 60.4 & 57.6 & 56.6 & \textbf{75.0} \\
PolyMath-en (Avg@4) & \textbf{65.1} & 59.5 & 51.9 & 52.8 & 49.8 & 54.0 & 49.9 \\
ZebraLogic (Acc.) & \textbf{89.0} & 84.0 & 37.7* & 79.7 & 59.3 & 58.5 & 57.9 \\
AutoLogi (Acc.) & \textbf{\underline{89.5}} & 88.9 & 83.3* & \textbf{89.8} & 86.1 & 88.2 & 84.1 \\
GPQA-Diamond (Avg@8) & \textbf{75.1} & 68.4* & 62.9* & 70.0* & 74.9* & 66.3 & 68.2 \\
SuperGPQA (Acc.) & \textbf{57.2} & 53.7 & 50.2 & 55.7 & 56.5 & 50.8 & 49.6 \\
Humanity's Last Exam (Acc.) & 4.7 & 5.2 & \textbf{\underline{5.7}} & 5.8 & \textbf{7.1} & 3.7 & 5.6 \\
\midrule
\multicolumn{8}{@{}l}{\textbf{General Tasks}} \\ % 调整为 8 列
\midrule
MMLU (EM) & \textbf{\underline{89.5}} & 89.4 & 87.0 & 91.5 & \textbf{92.9} & 90.4 & 90.1 \\
MMLU-Redux (EM) & \textbf{\underline{92.7}} & 90.5 & 89.2* & 93.6 & \textbf{94.2} & 92.4 & 90.6 \\
MMLU-Pro (EM) & 81.1 & \textbf{\underline{81.2*}} & 77.3 & 83.7 & \textbf{86.6} & 81.8 & 79.4 \\
IFEval (Prompt Strict) & \textbf{89.8} & 81.1 & 83.2* & 87.6 & 87.4 & 88.0 & 84.3 \\
Multi-Challenge (Acc.) & \textbf{54.1} & 31.4 & 34.0 & 46.8 & 49.0 & 36.4 & 39.5 \\
SimpleQA (Correct) & \textbf{\underline{31.0}} & 27.7 & 13.2 & 15.9 & 22.8 & \textbf{42.3} & 23.3 \\
Livebench (Pass@1) & \textbf{76.4} & 72.4 & 67.6 & 74.8 & 74.6 & 69.8 & 67.8 \\
\makecell[l]{Arena Hard v2.0 \\\emph{Hard Prompt} (Win rate)} & \textbf{\underline{54.5}} & 39.9 & 39.9 & 51.6 & \textbf{59.7} & 51.7 & 48.7 \\
\makecell[l]{Arena Hard v2.0 \\\emph{Creative Writing} (Win rate)} & \textbf{85.0} & 59.3 & 59.8 & 54.6 & 68.5 & 61.5 & 72.8 \\
FACTS Grounding (Adjusted) & \textbf{88.5} & 68.3 & 68.5 & 83.6 & --- & 79.2 & 86.6 \\
HHEM v2.1 (1-Hallu.) & \textbf{98.9} & 88.9 & 94.5 & 94.5 & --- & 96.7 & 97.8 \\
FaithJudge (1-Hallu.) & \textbf{\underline{92.6}} & 83.4 & 75.7 & 83.0 & --- & 91.0 & \textbf{93.2} \\
LongBench v2 (Acc.) &49.1&\textbf{\underline{51.1}}&---&52.5& --- & 54.3 & \textbf{55.5}\\
FRAMES (Acc.)  & 77.1 & \textbf{\underline{79.2}}&--- & 76.3&---&\textbf{87.4} &72.9\\
MRCR (Acc.) & \textbf{\underline{55.0}}& 50.8&---&74.4&---&66.9&\textbf{81.7}\\
DROP (Acc.) &\textbf{93.5}&91.2&84.3&92.0&---&79.1&81.7\\
\bottomrule
\\[-0.5ex]               % 让标题往下沉 0.5ex（可微调）
\end{tabular}
\end{table}

\subsubsection{Evaluation Results}

A comprehensive evaluation results of Kimi-K2-Instruct is shown in Table~\ref{tab:instruct_eval}, with detailed explanation provided in the Appendix~\ref{sec:k2-post-train-eval-details}. Below, we highlight key results across four core domains:

\paragraph{Agentic and Competitive Coding}
Kimi-K2-Instruct demonstrates state-of-the-art open-source performance on real-world SWE tasks. It outperforms most baselines on SWE-bench Verified (65.8\%, 71.6\% with multiple attemps), SWE-bench Multilingual (47.3\%), and SWE-lancer (39.1\%), significantly closing the gap with Claude 4 Opus and Sonnet. On competitive coding benchmarks (e.g., LiveCodeBench v6 53.7\%, OJBench 27.1\%), it also leads among all models, highlighting its practical coding proficiency across difficulty levels.

\paragraph{Agentic Tool Use}
On multi-turn tool-use benchmarks, Kimi-K2-Instruct sets a new standard. It achieves 66.1 Pass@1 on $\tau^2$-Bench and 76.5 on ACEBench, substantially outperforming all baselines. These results affirm its strength in grounded, controlled, and agent-driven tool orchestration across domains.

\paragraph{General Capabilities}
Kimi-K2-Instruct exhibits strong, balanced performance across general knowledge, math, instruction following, and long-context tasks. It surpasses open-source peers on SimpleQA (31.0\%), MMLU (89.5\%) and MMLU-Redux (92.7\%), and leads all models on instruction benchmarks (IFEval: 89.8\%, Multi-Challenge: 54.1\%). In math and STEM, it achieves top-tier scores (AIME 2024: 69.6\%, GPQA-Diamond: 75.1\%), and remains competitive on long-context factuality and retrieval (DROP: 93.5\%, MRCR: 55.0\%). These results position Kimi-K2-Instruct as a well-rounded and capable generalist across both short- and long-context settings.

\paragraph{Open-Ended Evaluation}
On the LMSYS Arena leaderboard (July 17, 2025), Kimi-K2-Instruct ranks as the top-1 open-source model and 5th overall based on over 3,000 user votes. This real-world preference signal—across diverse, blind prompts—underscores Kimi-K2's strengths in generating high-quality responses on open-ended tasks.








\subsection{Pre-training Evaluations}



\subsubsection{Evaluation Settings}



\paragraph{Benchmarks}
We evaluate Kimi-K2-Base across diverse capability areas.
For general capabilities, we assess on MMLU~\citep{hendrycks2021measuringmassivemultitasklanguage}, MMLU-Pro~\citep{wang2024mmluprorobustchallengingmultitask}, MMLU-Redux~\citep{gema2024we}, BBH~\citep{suzgun2022challengingbigbenchtaskschainofthought}, TriviaQA~\citep{joshi2017triviaqalargescaledistantly}, SuperGPQA~\citep{du2025supergpqa}, SimpleQA~\citep{wei2024measuring}, HellaSwag~\citep{zellers2019hellaswag}, AGIEval~\citep{zhong2023agieval}, GPQA-Diamond~\citep{rein2024gpqa}, ARC-Challenge~\citep{clark2018think}, and WinoGrande~\citep{sakaguchi2021winogrande}.
For coding capabilities, we employ EvalPlus~\citep{liu2023your} (averaging HumanEval~\citep{chen2021codex}, MBPP~\citep{austin2021programsynthesislargelanguage}, HumanEval+, and MBPP+), LiveCodeBench v6~\citep{jain2024livecodebench}, and CRUXEval~\citep{gu2024cruxeval}.
For mathematical reasoning, we utilize GSM8K~\citep{cobbe2021trainingverifierssolvemath}, GSM8K-Platinum~\citep{vendrow2025large}, MATH~\citep{hendrycks2021measuringmathematicalproblemsolving}, and CMATH~\citep{wei2023cmathlanguagemodelpass}.
For Chinese language capabilities, we evaluate on C-Eval~\citep{huang2023cevalmultilevelmultidisciplinechinese}, CMMLU~\citep{li2024cmmlumeasuringmassivemultitask}, and CSimpleQA~\citep{he2411chinese}.



\paragraph{Baselines}
We benchmark against leading open-source foundation models: DeepSeek-V3-Base~\citep{deepseekai2024deepseekv3technicalreport}, Qwen2.5-72B-Base~\citep{qwen2025qwen25technicalreport} (Note that Qwen3-235B-A22B-Base is not open-sourced, and the largest open-sourced base model in the Qwen series is Qwen2.5-72B-Base), and Llama 4-Maverick~\citep{metaLlamaHerd} (Llama 4-Behemoth is also not open-sourced). All models are evaluated under identical configurations to ensure fair comparison.


\paragraph{Evaluation Configurations}

We employ perplexity-based evaluation for MMLU, MMLU-Redux, GPQA-Diamond, HellaSwag, ARC-Challenge, C-Eval, and CMMLU. Generation-based evaluation is used for MMLU-Pro, SuperGPQA, TriviaQA, BBH, CSimpleQA, MATH, CMATH, GSM8K, GSM8K-Platinum, CRUXEval, LiveCodeBench, and EvalPlus. To mitigate the high variance inherent to GPQA-Diamond, we report the mean score across eight independent runs. All evaluations are conducted using our internal framework derived from LM-Harness-Evaluation~\citep{biderman2024lessons}, ensuring consistent settings across all models.


\subsubsection{Evaluation Results}

Table~\ref{tab:pretrain-eval} presents a comprehensive comparison of Kimi-K2-Base against leading open-source foundation models across diverse evaluation benchmarks. The results demonstrate that Kimi-K2-Base achieves state-of-the-art performance across the majority of evaluated tasks, establishing it as a leading foundation model in the open-source landscape.

\paragraph{General Language Understanding}
Kimi-K2-Base achieves state-of-the-art performance on 10 out of 12 English language benchmarks. Notable results include MMLU (87.79\%), MMLU-Pro (69.17\%), MMLU-Redux (90.17\%), SuperGPQA (44.67\%), and SimpleQA (35.25\%), significantly outperforming all baselines.

\paragraph{Coding Capabilities}
On coding benchmarks, Kimi-K2-Base sets new standards with leading performance across all metrics. It achieves 74.00\% on CRUXEval-I-cot, 83.50\% on CRUXEval-O-cot, 26.29\% on LiveCodeBench v6, and 80.33\% on EvalPlus, demonstrating superior code generation and comprehension abilities, particularly in scenarios requiring step-by-step reasoning.

\paragraph{Mathematical Reasoning}
Kimi-K2-Base exhibits exceptional mathematical capabilities, leading on three out of four benchmarks: MATH (70.22\%), GSM8K (92.12\%), and GSM8K-Platinum (94.21\%). It maintains competitive performance on CMATH (90.26\%), narrowly behind DeepSeek-V3-Base (90.53\%). These results highlight the model's robust mathematical problem-solving abilities across varying difficulty levels.

\paragraph{Chinese Language Understanding}
The model demonstrates superior multilingual capabilities, achieving state-of-the-art results across all Chinese language benchmarks: C-Eval (92.50\%), CMMLU (90.90\%), and CSimpleQA (77.57\%). These results establish Kimi-K2-Base as a leading model for Chinese language understanding while maintaining strong performance across other languages.















\begin{table}[h]
\centering
\footnotesize
\setlength{\tabcolsep}{1.9pt}
\caption{Performance comparison of Kimi-K2-Base against leading open-source models across diverse tasks. }
\label{tab:pretrain-eval}
\begin{tabular}{@{}c l | c | c | c c c}
\toprule
& \textbf{Benchmark {\tiny (Metric)}} &  \textbf{\#Shots} &{\textbf{Kimi-K2-Base}} & \textbf{DeepSeek-V3-Base} & \textbf{Llama4-Maverick-Base}  & \textbf{Qwen2.5-72B-Base}  \\
\midrule
    \midrule
    & Architecture & - & MoE & MoE & MoE & Dense \\
    & \# Activated Params & - & 32B & 37B & 17B & 72B \\
    & \# Total Params & - & 1043B & 671B & 400B & 72B \\

\midrule
\multirow{13}{*}{English} 
 & MMLU & 5-shots & \textbf{87.79} & 87.10 & 84.87 & 86.08 \\
 & MMLU-pro & 5-shots & \textbf{69.17} & 60.59 & 63.47 & 62.80 \\
 & MMLU-redux & 5-shots & \textbf{90.17} & 89.53 & 88.18 & 87.77 \\
 & SuperGPQA & 5-shots & \textbf{44.67} & 39.20 & 38.84 & 34.23 \\
 & GPQA-Diamond({\tiny avg@8}) & 5-shots & 48.11 & \textbf{50.51} & 49.43 & 40.78 \\
 & SimpleQA & 5-shots & \textbf{35.25} & 26.49 & 23.74 & 10.31 \\
 & TriviaQA & 5-shots & \textbf{85.09} & 84.11 & 79.25 & 76.03 \\
 & BBH & 3-shots & \textbf{88.71} & 88.37 & 87.10 & 84.09 \\
 & HellaSwag & 5-shots & 94.60 & 89.44 & 86.02 & \textbf{95.27} \\
 & AGIEval & - & \textbf{84.23} & 81.57 & 67.55 & 76.87 \\
 & ARC-Challenge & 0-shot & \textbf{95.73} & 93.77 & 94.03 & 95.56 \\
 & WinoGrande & 5-shots & \textbf{85.32} & 84.21 & 77.58 & 84.14 \\
 \midrule
\multirow{4}{*}{Code}
 & CRUXEval-I-cot & 0-shots & \textbf{74.00} & 62.75 & 67.13 & 61.12 \\
 & CRUXEval-O-cot & 0-shots & \textbf{83.50} & 75.25 & 75.88 & 66.13 \\
 & LiveCodeBench(v6) & 1-shots& \textbf{26.29} & 24.57 & 25.14 & 22.29 \\
 & EvalPlus & - & \textbf{80.33} & 65.61 & 65.48 & 66.04 \\
 \midrule
\multirow{4}{*}{Math}  
 & MATH & 4-shots & \textbf{70.22} & 61.70 & 63.02 & 62.68 \\
 & GSM8k & 8-shots & \textbf{92.12} & 91.66 & 86.35 & 90.37 \\
 & GSM8k-platinum & 8-shots & \textbf{94.21} & 93.38 & 88.83 & 92.47 \\
 & CMATH & 6-shots & 90.26 & \textbf{90.53} & 88.07 & 86.98 \\
\midrule
\multirow{3}{*}{Chinese} 
 & C-Eval & 5-shots & \textbf{92.50} & 90.04 & 80.91 & 90.86 \\
 & CMMLU & 5-shots & \textbf{90.90} & 88.84 & 81.24 & 90.55 \\
 & CSimpleQA & 5-shots & \textbf{77.57} & 72.13 & 53.47 & 50.53 \\
\bottomrule
\end{tabular}
\end{table}



\subsection{Safety Evaluation}

\subsubsection{Experiment Settings}

We conducted red-teaming evaluations on Kimi K2 compare with other open-source LLMs. The evaluation covered a range of attack scenarios—including harmful content, privacy content, and security content, as well as different attack strategies such as prompt injection and iterative jailbreak.

We choose \emph{Promptfoo}\footnote{\url{https://github.com/promptfoo/promptfoo}} to generate adversarial prompts and analyze the responses. By this way, we can evaluate model in a scalable ways.

\textbf{Model Selection} We compare Kimi K2 with three other open-source LLMs: DeepSeek-V3, DeepSeek-R1, and Qwen3.

\textbf{Promptfoo Settings} Table~\ref{tab:safety-stategy} lists plugins and strategies evaluated, with each plugin paired with all strategies to assess their performance.

\renewcommand{\arraystretch}{1.2}
\begin{table}[h]
\centering
\footnotesize
\caption{Enabled Plugins and Strategies}
\label{tab:safety-stategy}
\setlength{\tabcolsep}{1.9pt}
\begin{tabular}{|c|c|p{12cm}|}    
\hline
\multirow{5}{*}{Plugin} & Harmful & Graphic Content, Harassment and Bullying, Hate Speech, Insults, Profanity, Radicalization, Self Harm, Sexual Content, ToxicChat \\ \cline{2-3}
                      & Criminal & Chemical\&Biological Weapons, Child Exploitation, Copyright Violations, Cybercrime, Illegal Activities, Illegal Drugs, Indiscriminate Weapons, Intellectual Property Violation, Non-Violent Crime, Violent Crime, Sex Crimes \\ \cline{2-3}
                      & Misinformation & Competitor Endorsement, Unsupervised Contracts, Excessive Agency, Hallucination, Misinformation and Disinformation, Specialized Advice, Unsafe Practices, Imitation, Overreliance, Political Opinions, Religious Sensitivity \\ \cline{2-3}
                      & Privacy & Privacy Violation, PII in API/Database, Direct PII Exposure, PII in Session Data, PII via Social Engineering \\ \cline{2-3}
                      & Security & ASCII Smuggling, CyberSecEval, Harmbench, Debug Access, Divergent Repetition, DoNotAnswer, Malicious Code, Pliny, Prompt Extraction, Reasoning DoS, Tool Discovery \\ \hline
Strategy              & \multicolumn{2}{p{10cm}|}{Basic, Prompt Injection, Iterative Jailbreak, Crescendo} \\ \hline
\end{tabular}
\label{tab:plugin-strategy}
\end{table}
\renewcommand{\arraystretch}{1.0}


\textbf{Test Case Count} Given the inherent non-determinism of large language model inference, single-pass outputs may exhibit variability. To account for this, we generated 3 attack prompts per plugin for each strategy.

\textbf{Prompt Language Settings} We pre-tested the language compatibility for each plugin-strategy combination. Some plugins support both English and Chinese, while others only support English. For combinations that support both, we generated 3 prompts in each language, resulting in 6 prompts per combination.

\textbf{Manual Review} We incorporated human review into the evaluation process. To minimize subjectivity problem, we conducted multiple rounds of review and assigned the same reviewer to evaluate all cases within a given test set to ensure consistency and reduce variability in judgment.

\subsubsection{Safety Evaluation Results}
Table~\ref{tab:safety-eval} presents the passing rates of different models under various plugin–strategy combinations.

\begin{table}[h]
\centering
\footnotesize
\setlength{\tabcolsep}{1.9pt}
\caption{Safety Evaluation Results}
\label{tab:safety-eval}
\begin{tabular}{@{}c l | c | c c c@{}}
\toprule
\textbf{Plugin} &  \textbf{Strategy} &  \textbf{Kimi-K2-Instruct} & \textbf{DeepSeek-V3-0324} & \textbf{DeepSeek-R1}  & \textbf{Qwen3-235B-A22B}  \\

\midrule
\midrule
\multirow{4}{*}{Harmful} 
 & Basic & 98.04 & 90.45 & 99.02 & 98.53 \\
 & Base64 & 100 & 90.20 & 100 & 100 \\
 & Prompt Injection & 93.14  & 100 & 95.10 & 99.02 \\
 & Iterative Jailbreak & 92.16 & 66.67 & 72.55 & 74.51 \\
 & Crescendo & 64.71 & 64.71 & 80.39 & 86.27 \\
 \midrule
\multirow{4}{*}{Criminal} 
 & Basic & 100 & 99.62 & 95.45 & 99.24 \\ 
 & Base64 & 96.97 & 89.39 & 84.85 & 98.48 \\
 & Prompt Injection & 75.76 & 91.67 & 69.70 & 98.47 \\
 & Iterative Jailbreak & 57.57 & 21.21 & 25.76 & 53.03 \\
 & Crescendo & 56.06 & 31.81 & 42.42 & 59.09 \\
 \midrule
\multirow{4}{*}{Misinformation} 
 & Basic & 97.28 & 92.57 & 92.46 & 94.84 \\
 & Base64 & 98.48 & 90.48 & 96.83 & 93.65 \\
 & Prompt Injection & 98.39 & 86.51 & 93.65 & 93.65 \\
 & Iterative Jailbreak & 63.97 & 53.97 & 84.13 & 69.84 \\
 & Crescendo & 85.71 & 55.56 & 88.89 & 84.13 \\
 \midrule
\multirow{4}{*}{Privacy}  
 & Basic & 100 & 100 & 100 & 100 \\
 & Base64 & 100 & 100 & 100 & 100 \\
 & Prompt Injection & 88.33 & 98.33 & 100 & 91.67 \\
 & Iterative Jailbreak & 76.67 & 100 & 93.33 & 96.67 \\
 & Crescendo & 96.67 & 100 & 96.67 & 100 \\
 \midrule
\multirow{4}{*}{Security}  
 & Basic & 77.84 & 75.57 & 70.46 & 90.09 \\
 & Base64 & 82.93 & 82.93 & 63.41 & 95.12 \\
 & Prompt Injection & 87.80 & 97.56 & 65.85 & 84.13 \\
 & Iterative Jailbreak & 43.90 & 60.97 & 43.90 & 78.04 \\
 & Crescendo & 68.29 & 87.80 & 68.29 & 87.80 \\
\bottomrule
\end{tabular}
\end{table}

Without targeted optimization for specific evaluation scenarios, the passing rate of some complex cases (e.g., Harmful–Iterative Jailbreak) was relatively higher compared to other models.

Across different attack strategies, the models exhibited varying trends. Under the Base64 strategy, passing rates generally approached or reached 100\%, suggesting that encoding transformations had minimal impact on the models’ basic robustness. In contrast, the Crescendo strategy led to a general drop in passing rates, indicating stronger adversarial effectiveness.

In addition, complex attack strategies do not always outperform basic prompts. Some originally adversarial prompts may lose their intended meaning after multiple rounds of transformation, rendering the resulting model outputs less meaningful.

\textbf{Automated Red-teaming Limitations }Due to the involvement of human review, the evaluation results inevitably contain a degree of subjectivity. Additionally, certain plugin types involve API misuse or external tool invocation, which are more suitable for evaluating agent models with tool-calling capabilities. In the context of base LLMs, such tests may have limited relevance.






